% !TeX program = pdflatex
% !TeX encoding = UTF-8
% 根据 docs/template/main.tex 的版式与章节结构创建。
% 正式竞标书；中英文内容同步维护。
% 编译：使用 pdfLaTeX，连续编译两次更新目录。
% 成员资料来自 docs/character/；不得用往年成员信息填补未知字段。
\documentclass[a4paper,12pt]{article}
\usepackage[utf8]{inputenc}
\usepackage{amsmath}
\usepackage{array}
\usepackage[table]{xcolor}
\usepackage{needspace}
\usepackage{fancyhdr}
\usepackage{graphicx}
\usepackage[hidelinks]{hyperref}
\usepackage{tocloft}
\renewcommand{\labelitemi}{\ensuremath{\bullet}}
\setlength{\headheight}{14.5pt}
\addtolength{\headsep}{-2.5pt}
% Allow a little extra spacing on difficult lines without changing template margins.
\setlength{\emergencystretch}{3em}

% 统一使用本组名称，不沿用往年模板的组号。
\newcommand{\BidTeam}{UNNC GRP Team 10}
\newcommand{\BidDate}{October 15, 2026}
\newcommand{\BidHeaderDate}{October 15, 2026}
\newcommand{\PersonalInformation}[2]{%
  \subsection*{Personal Information}
  \if\relax\detokenize{#1}\relax\else
  \noindent\textbf{Student ID:} #1\\
  \fi
  \textbf{Email:} #2\par
}
\pagestyle{fancy}
\fancyhf{}
\fancyhead[L]{\BidTeam}
\fancyhead[R]{\BidHeaderDate}
\fancyfoot[C]{\thepage}
\renewcommand{\headrulewidth}{0.4pt}
\renewcommand{\footrulewidth}{0pt}

\newcommand{\BidProjectCode}{P21}
\newcommand{\BidProjectName}{AI-Driven Post-Investment Quarterly Report Analysis for Private Equity Portfolios}
\newcommand{\BidTitle}{Bid for [\BidProjectCode]}
\title{\BidTitle}
\author{\BidTeam}
\date{\BidDate}
\hypersetup{pdftitle={Bid for [P21]},pdfauthor={UNNC GRP Team 10}}

\begin{document}
\maketitle
\thispagestyle{plain}
\tableofcontents
\newpage

\section{Introduction}
\noindent\textbf{Project:} \BidProjectName
\par\medskip
Our interest in AI and finance is supported by experience directly relevant to P21's document-to-report workflow. Yun Xin NG has extracted data from PDFs and compared vision-language models for document understanding; Ziqian XIONG brings model experiments and Flask/React application development; Jinyao LU contributes LLM-related research and technical writing. Together with the team's annotation, database, and CI/CD skills, these experiences give UNNC GRP Team 10 a practical starting point for the project.

We want to turn complex quarterly information into clear, checkable support for post-investment analysis. The system will connect PDF/Word extraction and LLM-based structuring with company and industry analysis, peer comparisons, graded risks, and standardised report generation. Key figures and findings will link to source evidence for analyst review.

We will first demonstrate the complete workflow for one company across multiple quarters, then extend the analysis and reporting coverage. Ziqian XIONG will coordinate shared learning and early demonstrations, while members contribute their relevant experience and investigate unfamiliar methods together. Financial definitions, peer-selection criteria, and risk rules will be developed with supervisor or domain-reviewer input, and the system will be evaluated on authorised samples.

\medskip
\noindent\textbf{Recommended reading: our P21 bid plan website}\par
We have prepared a dedicated website for UNNC GRP Team 10's P21 bid plan. We recommend that reviewers visit the website for a convenient overview of our proposal alongside this report:\par
\noindent\url{https://team10-p21.vercel.app/}

\newpage
\section{Teammates}
P21 requires document extraction, dependable analysis, and a usable reporting application. Our seven-member team brings experience relevant to each stage, as outlined below. Led by Ziqian XIONG, we will coordinate this work through shared data definitions, weekly integration reviews during concentrated development, and early demonstrations; the low-workload period will use brief asynchronous check-ins. The following profiles present the experience supporting our project plan; responsibilities are set out in the work-distribution section.

\subsection{Ziqian XIONG}
\PersonalInformation{20721428}{\href{mailto:scyzx7@nottingham.edu.cn}{\nolinkurl{scyzx7@nottingham.edu.cn}}}
\noindent\textbf{Group role:} Leader\par

\subsection*{Introduction}
Ziqian XIONG is a Computer Science undergraduate at UNNC and our team leader. His medical image segmentation research involves model implementation, dataset checking, and comparative evaluation. His 4DGS tool and MomoFocus application add experience in Python/Flask backends and React/TypeScript interfaces. For P21, these experiences support comparing extraction methods and connecting model outputs with an interactive verification and reporting application. Personal website: \url{https://qqqqqprofile.vercel.app/}.

\subsection*{Internship Experience}
\begin{itemize}
  \item \textbf{Pony.ai algorithm internship (July--September 2026):} In the AI department, investigated VGGT-related 3D/4D reconstruction methods and compared multi-view geometry and depth estimation against an existing DA3/3D Gaussian Splatting pipeline. Completed feasibility validation and engineering integration, replacing the original reconstruction workflow. Also wrote documentation on agent-assisted development workflows.
\end{itemize}

\subsection*{Project Experience}
\begin{itemize}
  \item \textbf{4DGS-Edit-and-Compare (July--September 2026):} Independently developed a 4D Gaussian Splatting point-cloud editing and comparison tool. It supports rotation, translation, and deletion of selected regions, editing and comparing multiple point clouds, and quantitative evaluation. \textbf{Technology stack:} Python/Flask for the backend, JavaScript/Three.js/WebGL for the interactive frontend, NumPy/SciPy for numerical processing, PyTorch/plyfile for point-cloud file handling, and Docker for packaging. Developed independently during the Pony.ai internship, the tool is reported on his website as being used by the company. \href{https://github.com/qianqqqqqXZQ/4DGS-Edit-and-Compare}{Project repository}.
  \item \textbf{MomoFocus (September--October 2026):} Developed a personal productivity application integrating to-do lists, notes, and a Pomodoro timer. \textbf{Technology stack:} React, TypeScript, and CSS for the interface; Electron for the desktop shell; Vite for development and frontend builds; and electron-builder for desktop packaging. The implementation uses localStorage for local persistence and dnd-kit for sortable lists, providing relevant experience in interactive application development and desktop delivery. \href{https://github.com/qianqqqqqXZQ/MomoFocus}{Project repository}.
\end{itemize}

\subsection*{Research Experience}
\begin{itemize}
  \item \textbf{Medical image segmentation research:} Implemented and debugged CNN--Transformer components; prepared and checked datasets; evaluated Dice, IoU, and inference speed. \textbf{Technology stack:} Python, PyTorch, and OpenMMLab's MMSegmentation, MMCV, and MMEngine. A related paper is under submission to CVIU with him as second author.
  \item \textbf{Dynamic dual-branch neural network research:} Investigating a multi-task network for keyword spotting and speaker verification, with early-exit scheduling and FPGA resource reallocation. Network design, scheduling, and resource-allocation strategies remain under evaluation.
\end{itemize}

\subsection*{Skill Set}
\begin{itemize}
  \item \textbf{Programming languages:} Python, C/C++, Java, Haskell, JavaScript, TypeScript, Lean, and Verilog.
  \item \textbf{Databases:} MySQL.
  \item \textbf{Research and implementation:} PyTorch; OpenMMLab (MMSegmentation, MMCV, and MMEngine); CNN--Transformer models; data preprocessing and comparative model evaluation; Retrieval-Augmented Generation (RAG).
  \item \textbf{Application development:} HTML, CSS, React, Electron, Vite, Flask, Three.js/WebGL, NumPy, and SciPy.
  \item \textbf{Engineering and documentation:} Docker packaging, electron-builder, and technical documentation for agent-assisted development workflows, including harness and loop engineering.
\end{itemize}

\subsection{Yun Xin NG}
\PersonalInformation{20709091}{efyyn8@nottingham.edu.cn}
\noindent\textbf{Group role:} Member\par

\subsection*{Introduction}
Yun Xin NG is a second-year Computer Science student on a one-year exchange at UNNC after completing her first year at the University of Nottingham Malaysia. At YTL AI Labs, she extracted data from websites and PDFs and used existing benchmarks to compare open-source vision-language models for document understanding and text extraction. This experience directly supports P21's document-parsing and model-selection work. Her machine learning projects and MySQL-backed course application also provide a foundation for preparing evaluation samples and storing extracted records.

\subsection*{Internship Experience}
\begin{itemize}
  \item \textbf{YTL AI Labs data-engineering internship:} Extracted data from web pages and PDF documents, shortlisted open-source vision-language models against application needs, and compared document-understanding and text-extraction capabilities using existing benchmarks. Selected a model by considering both benchmark performance and practical suitability.
\end{itemize}

\subsection*{Project Experience}
\begin{itemize}
  \item \textbf{Machine learning project:} Prepared training/test splits and trained models, gaining experience in dataset preparation and the experimental workflow.
  \item \textbf{Internship Management System:} Developed a first-year course application using HTML, CSS, Java, and MySQL, connecting page development with application logic and persistent data.
\end{itemize}

\subsection*{Skill Set}
\begin{itemize}
  \item \textbf{Programming languages:} Python, Java, and C/C++.
  \item \textbf{Web and database technologies:} HTML, CSS, and MySQL.
  \item \textbf{Data engineering:} Data extraction from web pages and PDFs, dataset preparation, training/test splitting, and Retrieval-Augmented Generation (RAG).
  \item \textbf{Models and AI technologies:} Model training and benchmark-based vision-language model evaluation and selection.
\end{itemize}

\subsection{Jinyao LU}
\PersonalInformation{}{scyl@nottingham.edu.cn}
\noindent\textbf{Group role:} Member\par

\subsection*{Introduction}
Jinyao LU is a Computer Science and Artificial Intelligence undergraduate at UNNC with LLM-related research and academic-writing experience. For P21, this background provides a basis for investigating evidence-grounded report generation, translating agreed reporting requirements into template and review criteria, and communicating evaluation results. His Python/NumPy and software-engineering foundations support participation in implementation and testing.

\subsection*{Internship Experience}
\begin{itemize}
  \item \textbf{R\&D internship at Zhejiang Dishi Information Technology (July--August 2024):} Learned the software-development workflow and platform business logic; investigated VR/AR applications in the water sector.
  \item \textbf{Internship at Zhejiang Zhongkong Information Industry (June--August 2025):} Monitored train status, participated in routine maintenance of an information-publishing platform and operational data recording, and learned to operate the dispatch system.
\end{itemize}

\subsection*{Project Experience}
\begin{itemize}
  \item \textbf{Douyin AI hackathon organisation (March--April 2026):} Served on the campus core organising team, coordinating participating teams and event processes.
\end{itemize}

\subsection*{Research Experience}
\begin{itemize}
  \item \textbf{AI research:} First author of a paper on LLMs and autonomous driving accepted by an international journal.
\end{itemize}

\subsection*{Skill Set}
\begin{itemize}
  \item \textbf{Programming languages:} Python, Java, and C/C++.
  \item \textbf{Databases:} MySQL.
  \item \textbf{Technologies and foundations:} NumPy data processing, basic web development, data structures, algorithms, and software engineering.
  \item \textbf{Development and reporting tools:} Environment setup and basic use of C and Java programming tools; Excel charts and PowerPoint presentations.
  \item \textbf{Technical communication:} Academic writing and presentation of research findings.
\end{itemize}

\subsection{Oswald Yiwen CHONG}
\PersonalInformation{20792699}{\href{mailto:hcyoc1@nottingham.edu.cn}{\nolinkurl{hcyoc1@nottingham.edu.cn}}}
\noindent\textbf{Group role:} Member\par

\subsection*{Introduction}
Oswald Yiwen CHONG is from Malaysia and studies Computer Science with Artificial Intelligence at the University of Nottingham. He is passionate about building practical software and keen to keep learning through technical challenges and teamwork. His blockchain-based bouldering social application provides software-development experience, complemented by GitHub collaboration and CI/CD skills. For P21, these capabilities support connecting report upload, processing status, analysis, and report export into a tested application, with repeatable integration checks and deployment.

\subsection*{Project Experience}
\begin{itemize}
  \item \textbf{Blockchain-based bouldering social application:} Developed a Strava-like social application for bouldering, using blockchain to award user badges and incorporating NFT-related functionality. This provided practical experience in applying decentralised technology to a user-facing application.
\end{itemize}

\subsection*{Skill Set}
\begin{itemize}
  \item \textbf{Programming languages:} Python, Java, and C/C++.
  \item \textbf{Databases:} MySQL.
  \item \textbf{Application technologies:} Application development, blockchain, and decentralised technologies.
  \item \textbf{Version control and collaboration:} Proficient use of GitHub for version control and team collaboration.
  \item \textbf{Continuous integration and continuous deployment (CI/CD):} Able to set up pipelines that support code reliability and rapid software delivery.
\end{itemize}

\subsection{Ziqian ZHANG}
\PersonalInformation{20721592}{scyzz23@nottingham.edu.cn}
\noindent\textbf{Group role:} Member\par

\subsection*{Introduction}
Ziqian ZHANG is a third-year Computer Science student whose research involved image annotation and data preparation, with skills in Labelme and JSON annotation/data management. P21 likewise needs carefully checked reference data to assess extraction quality. We will adapt her annotation and data-organisation experience to report-field and source annotations, with a second reviewer checking financial values, units, and reporting periods.

\subsection*{Research Experience}
\begin{itemize}
  \item \textbf{AI-driven Analysis and Optimization of Performance on Magnets:} Participated in this research project, contributing image annotation and data preparation. Alongside this experience, her skills in Labelme and JSON annotation/data management provide a foundation for organising image labels and structured annotation records. These capabilities can support the preparation of labelled examples and evaluation materials, with attention to consistency between source images and their annotations.
\end{itemize}

\subsection*{Skill Set}
\begin{itemize}
  \item \textbf{Programming languages:} Python, Java, and C/C++ at a foundational level.
  \item \textbf{Databases:} MySQL.
  \item \textbf{Data tools and processing:} Labelme, JSON annotation and data management, image annotation, and dataset preparation.
\end{itemize}

\subsection{Yue WU}
\PersonalInformation{20717823}{\href{mailto:sgyyw10@nottingham.edu.cn}{\nolinkurl{sgyyw10@nottingham.edu.cn}}}
\noindent\textbf{Group role:} Member\par

\subsection*{Introduction}
Yue WU studies Computer Science at UNNC. Her CCN research and paper writing involve literature review and technical analysis, complemented by Python, Java, and SQLite skills. For P21, this combination provides a basis for researching analytical methods, implementing reviewed quarterly calculations and risk rules, and documenting their assumptions so that results can be checked.

\subsection*{Project Experience}
\begin{itemize}
  \item \textbf{Coursework development:} Applied programming and algorithmic problem-solving in software-related course projects, building experience in translating computing concepts into implementations.
\end{itemize}

\subsection*{Research Experience}
\begin{itemize}
  \item \textbf{CCN research and technical writing:} Worked on research and paper writing concerning \emph{The softwarization and interoperability of CCN}, developing literature-review, technical-analysis, and academic-writing experience.
\end{itemize}

\subsection*{Skill Set}
\begin{itemize}
  \item \textbf{Programming languages:} Java, Python, Haskell, C/C++, and JavaScript.
  \item \textbf{Web and database technologies:} HTML, CSS, MySQL, and SQLite.
  \item \textbf{Development and analysis tools:} IntelliJ IDEA, Anaconda, Visual Paradigm, LaTeX, and MATLAB.
  \item \textbf{Technical analysis:} Algorithm analysis, technical literature review, and academic writing.
\end{itemize}

\subsection{Xinyun HU}
\PersonalInformation{20720227}{scyxh8@nottingham.edu.cn}
\noindent\textbf{Group role:} Member\par

\subsection*{Introduction}
Xinyun HU is a third-year Computer Science student with programming, SQL, SQLite, and web-technology foundations, alongside an AI-enabled Connect Four project. For P21, these skills provide a basis for building forms that connect extracted indicators to database records and for presenting quarterly comparisons, missing data, and review status to analysts.

\subsection*{Project Experience}
\begin{itemize}
  \item \textbf{Connect Four Game with AI (July 2025):} The supplied project background includes an AI-enabled Connect Four game, providing a relevant setting for discussing the interaction between algorithmic behaviour and an interactive application.
\end{itemize}

\subsection*{Skill Set}
\begin{itemize}
  \item \textbf{Programming and query languages:} Java, Python, C/C++, Haskell, SQL, MATLAB, and MIPS Assembly.
  \item \textbf{Web and database technologies:} HTML, CSS, MySQL, and SQLite.
  \item \textbf{Development tools:} Git, VS Code, Eclipse, and Visual Paradigm.
\end{itemize}

\subsection{Why Our Team Fits This Project}
\begin{itemize}
  \item \textbf{Strong motivation to connect AI with finance:} Our shared interest in AI and financial analysis motivates us to engage closely with the business questions behind the software. We will turn this interest into regular discussion of authorised report examples, shared learning on financial indicators and post-investment workflows, and iterative refinement following supervisor or domain-reviewer feedback.
  \item \textbf{PDF extraction and model selection:} Yun Xin NG's PDF data extraction and benchmark-based document-model comparisons are directly relevant to the first stage of P21. Together with Ziqian XIONG's model implementation and comparative experiments, they provide a basis for testing parsing and extraction candidates on authorised quarterly reports.
  \item \textbf{Checkable analysis and evaluation:} Ziqian ZHANG's annotation and JSON data-preparation experience can support field-and-source reference samples. Yue WU's programming and technical-analysis background can support implementing quarterly comparisons and risk rules against domain-reviewed examples. This connects data-quality checks with reproducible analytical results.
  \item \textbf{Structured storage and a complete reporting application:} Yun Xin NG's MySQL-backed management system, Ziqian XIONG's Flask and React/TypeScript projects, Xinyun HU's SQL and web foundations, and Oswald Yiwen CHONG's application-development and CI/CD skills support the database, verification interface, integration, and deployment. Jinyao LU's LLM-related research and writing provide a basis for investigating report-generation methods and documenting template and review criteria.
  \item \textbf{A defined route to financial-domain validation:} We will agree indicator definitions, reporting bases, peer criteria, and risk rules with the supervisor or domain reviewer, then implement them as checked examples, calculations, and evidence-linked outputs. This gives our existing software and AI experience a concrete route into the project's investment-analysis requirements.
\end{itemize}

\newpage

\section{Ideas \& Design}
\subsection{Project Proposal}
We will develop an LLM-based web application that imports portfolio companies'
quarterly reports, extracts PDF/Word text and tables, and structures free
text, financial indicators, and business descriptions into a database.
The system will analyse company and industry quarterly changes, compare
peers, grade risks, and automatically generate standardised post-investment
reports using an LLM. Analysts will be able to verify extracted content
and review generated reports; key figures and risk findings will link to
source evidence.

The first version will connect report import, text/table extraction,
LLM-based structuring, verification, database storage,
quarter-on-quarter/year-on-year calculation, and automatic standardised
report generation for one company across multiple quarters. We will then
integrate anomaly attribution and industry quarterly environment changes,
including policy, market sentiment, and sector shifts, to assess benefits
or pressures. Peer matching by sector and stage will show relative
positions and quarterly ranking movements, alongside graded operational, compliance,
team-change, and performance-bet (VAM) risks. Adaptable report templates
will include hold/exit/impairment indications and value-creation suggestions
for analyst review.

\subsection{Technology Stack \& System Design}
% 技术实现通过报告样本、算例与领域审核进行验证。
\begin{enumerate}
  \item \textbf{Back-end Development}\par
  Python with FastAPI will expose import, verification, analysis, report
  generation, and export APIs. Background parsing and model tasks will return IDs for
  progress tracking.

  \noindent\textbf{Storage and provenance.}
  MySQL will hold companies, reporting periods, structured free-text
  information, financial indicators, business descriptions, evidence
  references, risk rules, and verification/review/version records. Original
  PDF/Word files and generated reports will use file or object storage;
  parsed text and table structures can initially use MySQL JSON fields.
  A separate MongoDB service is unnecessary initially. Indicators record
  original and normalised values,
  currency, unit, quarterly/cumulative and consolidated/parent-only basis,
  source location, verification status, and version. MySQL
  \texttt{DECIMAL} and Python \texttt{Decimal} will preserve monetary
  precision under explicit rounding rules.

  \noindent\textbf{Parsing and extraction.}
  PyMuPDF/pdfplumber will be evaluated for PDF text and tables, with
  python-docx for Word text and tables. PaddleOCR will process scanned pages
  and image text as needed; direct extraction takes priority on readable
  pages. Table JSON will retain headers, cell positions, units, and source
  locations. CSV will support import/export rather than replace this
  richer representation. An LLM will structure free text, financial
  indicators, and business descriptions from the extracted text and tables
  into a defined database schema, supported by parsing and validation rules.
  Unit, period, and source checks will route uncertain results to human
  verification.

  \noindent\textbf{Analysis and report generation.}
  Python calculations and configurable rules will handle financial
  comparisons, peer positions, and risk grading. Calculations will use
  compatible reporting bases and explicitly handle missing periods and
  zero or negative baselines. Industry and peer data must retain sources
  and publication dates. The LLM will automatically assemble company and
  industry analysis, peer benchmarking, graded risks, hold/exit/impairment
  indications, and value-creation suggestions into a standardised
  post-investment report template, using verified database records and
  supporting evidence. Explanations will separate reported causes from
  hypotheses requiring review. Numeric and citation checks precede
  analyst approval; corrections create new versions and invalidate affected
  analyses and drafts.

  \noindent\textbf{Optional semantic retrieval.}
  BGE-M3 will support Chinese/English evidence retrieval, evaluated on
  report samples for quality and resource use. FAISS will index vectors, with chunk IDs,
  sources, and versions in MySQL. Before ranking, restrict evidence by
  company, period, access rights, and valid version. Embeddings locate
  passages; financial extraction and calculation remain separate. Keep
  index updates/backups consistent with document versions; consider Qdrant
  if filtering or concurrency needs increase.

  \item \textbf{Front-end Development}\par
  React will provide document upload, task progress, side-by-side source
  verification, trend charts, peer comparisons, risk evidence, and report
  review/export. FastAPI responses will expose source references, missing
  data, and verification status. Unapproved reports will remain visibly
  marked as drafts.
\end{enumerate}

\noindent\textbf{Development tools and deployment:}
Use JetBrains DataGrip for database work, PyCharm or VS Code for development,
and Docker Compose for reproducible deployment. Select local/API LLMs by
extraction quality, cost, resources, and confidential-data handling requirements.

\medskip
\noindent\textbf{Core data flow:}
Import/store quarterly reports $\rightarrow$ text/table extraction (OCR as needed)
$\rightarrow$ LLM-based structuring of free text, financial indicators,
and business descriptions $\rightarrow$ verification and MySQL storage
$\rightarrow$ company/industry analysis, peer benchmarking, and risk grading
$\rightarrow$ LLM-generated standardised post-investment report
$\rightarrow$ numeric/source checks and approval $\rightarrow$ export.
Optional vector retrieval supplies supporting passages.

\medskip
\noindent\textbf{Domain validation and evaluation:}
During requirements analysis, we will establish authorised samples, indicator definitions, templates, industry/peer data, VAM targets, resource budgets, and validation criteria with the supervisor.
Report missing evidence explicitly. Evaluate extraction/source accuracy,
calculation correctness, risk errors, and report consistency. Add semantic
retrieval after the core workflow is stable.

\newpage
\section{Timeline \& Workflow and Work Distribution}
\subsection{Timeline \& Workflow}
Ziqian XIONG will lead weekly progress and integration reviews during concentrated development and use brief asynchronous check-ins in the low-workload period. Paired discussions of methods and code, and early demonstrations, will support supervisor feedback. The leader will help members break down unfamiliar tasks and resolve dependencies. Each member will prepare questions, report blockers early, and record decisions and follow-up actions.

The project runs from 21 October 2026 to final submission on 1 April 2027, with work organised by date ranges, workload, and demonstrable outputs. We use 16 December 2026--5 January 2027 to review documents and finish earlier tasks as needed. As January and February may overlap the winter vacation, 6 January--28 February is a low-workload period: at most one small asynchronous task per member per week, adjusted to availability, with no critical integration, required training runs, or participant trials. University holiday dates remain subject to formal notices. We bring the first demonstrable workflow and pilot experiments forward to 15 December, complete integration and evaluation in March, and reserve 31 March--1 April for submission checks.

Course milestones follow the supplied tentative key-dates table, subject to formal course notices.

\begin{figure}[!ht]
\centering
{\small
\setlength{\tabcolsep}{4pt}
\renewcommand{\arraystretch}{1.3}
\resizebox{\textwidth}{!}{%
\begin{tabular}{|l|*{7}{c|}}
\hline
\textbf{Phase} & \multicolumn{3}{c|}{2026} & \multicolumn{4}{c|}{2027} \\
\hline
Planning window & Oct & Nov & Dec & Jan & Feb & Mar & Apr \\
\hline
1. Requirements / sample audit & \cellcolor{cyan!45}21+ & \cellcolor{cyan!45}1--3 & \rule{0pt}{1.6ex} & \rule{0pt}{1.6ex} & \rule{0pt}{1.6ex} & \rule{0pt}{1.6ex} & \rule{0pt}{1.6ex} \\
2. Parsing / source verification & \rule{0pt}{1.6ex} & \cellcolor{cyan!45}4--24 & \rule{0pt}{1.6ex} & \rule{0pt}{1.6ex} & \rule{0pt}{1.6ex} & \rule{0pt}{1.6ex} & \rule{0pt}{1.6ex} \\
3. First report / rule pilots & \rule{0pt}{1.6ex} & \cellcolor{cyan!45}25+ & \cellcolor{cyan!45}1--15 & \rule{0pt}{1.6ex} & \rule{0pt}{1.6ex} & \rule{0pt}{1.6ex} & \rule{0pt}{1.6ex} \\
Review documents / finish earlier tasks & \rule{0pt}{1.6ex} & \rule{0pt}{1.6ex} & \cellcolor{gray!25}16+ & \cellcolor{gray!25}1--5 & \rule{0pt}{1.6ex} & \rule{0pt}{1.6ex} & \rule{0pt}{1.6ex} \\
4. Light review / documentation & \rule{0pt}{1.6ex} & \rule{0pt}{1.6ex} & \rule{0pt}{1.6ex} & \cellcolor{yellow!35}6+ & \cellcolor{yellow!35}1+ & \rule{0pt}{1.6ex} & \rule{0pt}{1.6ex} \\
5. Industry / risk integration & \rule{0pt}{1.6ex} & \rule{0pt}{1.6ex} & \rule{0pt}{1.6ex} & \rule{0pt}{1.6ex} & \rule{0pt}{1.6ex} & \cellcolor{cyan!45}1--9 & \rule{0pt}{1.6ex} \\
6. Accuracy / domain review & \rule{0pt}{1.6ex} & \rule{0pt}{1.6ex} & \rule{0pt}{1.6ex} & \rule{0pt}{1.6ex} & \rule{0pt}{1.6ex} & \cellcolor{cyan!45}10--16 & \rule{0pt}{1.6ex} \\
7. Release / report / demo & \rule{0pt}{1.6ex} & \rule{0pt}{1.6ex} & \rule{0pt}{1.6ex} & \rule{0pt}{1.6ex} & \rule{0pt}{1.6ex} & \cellcolor{cyan!45}17--30 & \rule{0pt}{1.6ex} \\
Submission checks & \rule{0pt}{1.6ex} & \rule{0pt}{1.6ex} & \rule{0pt}{1.6ex} & \rule{0pt}{1.6ex} & \rule{0pt}{1.6ex} & \cellcolor{cyan!45}31 & \cellcolor{cyan!45}1 \\
\hline
\end{tabular}}
}
\caption{Project timetable (dates within each month; \texttt{+} means through month end; grey: document review / earlier tasks; amber: low workload)}
\label{fig:timetable}
\end{figure}

\begin{itemize}
  \Needspace{3\baselineskip}
  \item \textbf{21 October -- 3 November 2026}
    \begin{itemize}
      \item Agree analyst journeys, report templates, indicator definitions, and acceptance criteria. Confirm access to multi-quarter company reports and industry, peer, and VAM evidence; parse a sample and record its source locations.
      \item Profile parsing and LLM resource needs and prepare compute/storage requests from 21 October. Submit ethics forms by 28 October, covering confidential reports, model/API data handling, consent, and planned analyst evaluation; begin participant work only after the required approval.
    \end{itemize}

  \Needspace{3\baselineskip}
  \item \textbf{4--24 November 2026}
    \begin{itemize}
      \item Publish the team project website by 5 November, using project descriptions and non-confidential examples. Design the MySQL indicator/evidence schema, FastAPI contracts, and verification interface.
      \item Implement PDF/Word text and table extraction, OCR for scanned pages, and LLM-based structuring with unit/period validation. Deliver a parsing baseline and checked reference set by 24 November; freeze the first indicator definitions. Prepare interface skeletons and reusable evaluation fixtures in parallel with this baseline.
    \end{itemize}

  \Needspace{3\baselineskip}
  \item \textbf{25 November -- 15 December 2026}
    \begin{itemize}
      \item Implement QoQ/YoY calculations on comparable reporting bases, handle missing periods, and add initial anomaly rules. By 15 December, demonstrate upload, verification, storage, calculation, and LLM report generation/export for one company across multiple quarters. Check that correcting a source figure updates the affected analysis and report.
      \item In parallel, prepare sourced industry/peer examples, pilot all four risk categories, and build chart and report-approval interface skeletons with reusable comparison scripts. Submit the interim report by 10 December, using evidence available then; archive document versions, pilot rules, and unresolved issues by 15 December.
    \end{itemize}

  \Needspace{3\baselineskip}
  \item \textbf{16 December 2026 -- 5 January 2027}
    \begin{itemize}
      \item Use this time to review documents, collect feedback, and finish earlier tasks as needed; no critical work or participant trials are scheduled. Low-workload activity continues from 6 January.
    \end{itemize}

  \Needspace{3\baselineskip}
  \item \textbf{6 January -- 28 February 2027}
    \begin{itemize}
      \item Review existing indicator definitions, source references, industry/peer metadata, and saved risk examples in small independent tasks. Draft the template, rule-library, and user documentation without adding new analysis modules.
      \item Use asynchronous check-ins to maintain extraction-error records and prepare test references, review instructions, and March priorities from existing materials. Make no full-application delivery commitment in January or February; resume concentrated integration on 1 March.
    \end{itemize}

  \Needspace{3\baselineskip}
  \item \textbf{1--9 March 2027}
    \begin{itemize}
      \item Connect the piloted industry analysis and sector/stage peer comparisons, showing relative positions and quarterly ranking changes. Add the tested operational, compliance, team-change, and VAM risk rules, with source evidence and clear missing-data states.
      \item By 9 March, connect upload/progress, source verification, charts, risk review, and report approval/export, including reviewed hold/exit/impairment indications and value-creation suggestions. Complete the core rule library; defer optional retrieval and valuation extensions.
    \end{itemize}

  \Needspace{3\baselineskip}
  \item \textbf{10--16 March 2027}
    \begin{itemize}
      \item Use the prepared held-out companies/periods and references to compare parsing and LLM extraction. Check source accuracy, units, quarterly/cumulative bases, calculations, risk false positives/negatives, and numeric/citation consistency.
      \item Run approved analyst or domain review and reliability tests in parallel. Check access isolation, updates after corrections, missing evidence, and task recovery; complete the evaluation report and essential fixes by 16 March.
    \end{itemize}

  \Needspace{3\baselineskip}
  \item \textbf{17--30 March 2027}
    \begin{itemize}
      \item Freeze features and rules on 17 March, regress import-to-export journeys and risk cases, and verify Docker Compose deployment in a fresh environment by 23 March. Complete template, indicator, rule, access, backup, and recovery documentation.
      \item Use 24--30 March to rehearse source verification, quarterly/peer changes, graded risks, and reviewed report export with authorised or anonymised examples. Complete the software, final report, adaptable templates, rule library, and manuals by 30 March.
    \end{itemize}

  \Needspace{3\baselineskip}
  \item \textbf{31 March -- 1 April 2027}
    \begin{itemize}
      \item Verify release/report/documentation consistency, filenames, access, and receipts; submit the team final report and software by 1 April. Prepare demonstration, presentation, and promotional assets for later course submissions.
    \end{itemize}

\end{itemize}

\clearpage

\subsection{Work Distribution}
We will divide the work according to our relevant experience and interests, help one another learn unfamiliar methods, and adjust tasks through progress discussions and supervisor feedback.
\begin{itemize}
  \item \textbf{Ziqian XIONG:} Coordinate the team and supervisor discussions; guide system design, compare extraction and LLM methods, and connect structured data with report generation.
  \item \textbf{Yun Xin NG:} Develop PDF/Word text and table extraction, compare OCR and document models, and retain source locations, drawing on her document-extraction and model-evaluation experience.
  \item \textbf{Jinyao LU:} Study reporting requirements and templates, contribute to report generation, and prepare review examples and documentation, drawing on his LLM research and writing experience.
  \item \textbf{Oswald Yiwen CHONG:} Develop backend services, database storage, and report export; connect application modules and support automated checks and deployment.
  \item \textbf{Ziqian ZHANG:} Prepare checked field and source samples, organise annotations, and record extraction and report errors, building on her research data-preparation experience.
  \item \textbf{Yue WU:} Implement agreed quarterly calculations, industry and peer comparisons, and risk rules; check examples and explain the methods with supervisor or domain-reviewer guidance.
  \item \textbf{Xinyun HU:} Build upload, source-verification, comparison, and report-review pages; connect them to backend data and show missing information and review status.
\end{itemize}

Yun Xin NG and Ziqian ZHANG will check extraction samples together; Ziqian XIONG and Jinyao LU will connect analysis with report templates. Financial definitions and risk interpretations will be checked with the supervisor or a domain reviewer.

\clearpage
\section{Conclusion}
Our strong interest in AI and finance gives us a clear motivation to pursue P21 and learn the domain in depth. Our team brings experience that aligns with P21's document-to-report workflow: PDF extraction and document-model evaluation, model implementation and experiments, annotation and structured data preparation, and database-backed application development with CI/CD. We will combine these capabilities to deliver the parsing and report-generation prototype, adaptable report template, and quarterly monitoring indicators and risk-rule base required by the project.

We will first demonstrate report import, structured storage, verification, and report generation for one company, then extend the company, industry, benchmarking, and risk analysis. Financial definitions and investment-related interpretations will be validated through supervisor or domain-reviewer feedback. Regular integration reviews and early demonstrations will turn our relevant experience into concrete results that can be assessed throughout development.

\end{document}

